% ============================================================
% CAPITOLO 36 - PROGRESSIONI MODULANTI
% ============================================================
\chapter{Progressioni modulanti}
\label{chap:progressioni_modulanti}

Una progressione, ripetendo un modello su gradi diversi, può portare da una tonalità a un'altra. Le progressioni modulanti uniscono i due strumenti dei capitoli precedenti: la sequenza (Capitolo \ref{chap:progressioni_armoniche}) e il pivot (Capitolo \ref{chap:modulazione}).

\section{Progressioni che modulano}
\label{sec:progressioni_modulanti}

\begin{definizione}[Progressione modulante]
	\label{def:progressione_modulante}
	Una \textbf{progressione modulante} è una sequenza armonica che, attraverso le ripetizioni trasposte del modello, conduce da una tonalità a un'altra. La modulazione avviene spesso in modo graduale, quasi impercettibile, grazie alla coerenza del disegno ripetuto.
\end{definizione}

\begin{regola}[Modulazione attraverso progressioni]
	\label{reg:modulazione_progressioni}
	Per modulare con una progressione:
	\begin{enumerate}
		\item si inizia la progressione nella tonalità di partenza;
		\item si continua il disegno introducendo gradualmente le alterazioni che appartengono alla nuova tonalità;
		\item si conferma la nuova tonalità con una cadenza autentica.
	\end{enumerate}
	Le progressioni per quinte discendenti sono particolarmente adatte, perché:
	\begin{itemize}
		\item ogni quinta può diventare la dominante della tonalità successiva;
		\item le alterazioni che creano le sensibili compaiono in modo naturale;
		\item il movimento è direzionale e porta alla nuova tonica.
	\end{itemize}
\end{regola}

Le sequenze reali dei moduli del Capitolo \ref{chap:moduli}, che ripetono il modello identico trasposto di un tono per volta, sono progressioni modulanti in senso stretto: il centro tonale cambia a ogni modulo.

\begin{esempio}[Progressione modulante da Do maggiore a Sol maggiore]
	\label{ex:progressione_modulante}
	Il modello VI-II-V-I si presenta prima in Do maggiore (La minore, Re minore, Sol, Do) e poi, una quinta più in alto, in Sol maggiore (Mi minore, La minore, Re, Sol). Il Do finale del primo modulo è anche il pivot: tonica di Do e IV grado di Sol. Il Fa$\sharp$ compare solo nella dominante di Sol (Re maggiore), e la cadenza V-I in Sol conferma la nuova tonalità.
	\musica{es_cap36_01}
\end{esempio}
